\documentclass{article}

\usepackage[a-2u]{pdfx}
\usepackage[utf8]{inputenc}
\usepackage[T1]{fontenc}
\usepackage{lmodern}
\usepackage[czech]{babel}
\usepackage{array}
\usepackage[ddmmyyyy]{datetime}
\usepackage{geometry}
\usepackage{enumitem}
\usepackage{url}
\usepackage{tikz}
\usetikzlibrary{positioning,fit}

% --- Makety obrazovek (kapitola Uživatelské rozhraní): styly a pomocná makra ---
\tikzset{
  wf frame/.style={draw=black!70, rounded corners=3pt, fill=white},
  wf bar/.style={fill=gray!20},
  wf card/.style={draw=black!45, rounded corners=2pt, fill=gray!6},
  wf pre/.style={draw=black!25, fill=white},
  wf field/.style={draw=black!55, fill=white, rounded corners=1pt},
  wf popup/.style={draw=black!55, fill=white},
  wf h/.style={font=\linespread{1}\fontsize{8}{9}\selectfont\sffamily\bfseries, inner sep=1pt, anchor=west},
  wf t/.style={font=\linespread{1}\fontsize{7}{8}\selectfont\sffamily, inner sep=1pt, anchor=west},
  wf tt/.style={font=\linespread{1}\fontsize{6}{7}\selectfont\sffamily, inner sep=1pt, anchor=west},
  wf err/.style={wf tt, text=red!65!black},
  wf code/.style={font=\linespread{1}\fontsize{5.5}{6.6}\selectfont\ttfamily, inner sep=1pt, anchor=north west, align=left},
  wf tab/.style={font=\linespread{1}\fontsize{6}{7}\selectfont\sffamily, inner sep=1pt, anchor=north west},
  wf btn/.style={draw=black!55, rounded corners=2pt, fill=gray!22, font=\linespread{1}\fontsize{6}{7}\selectfont\sffamily\bfseries, inner sep=2.5pt, minimum height=4.4mm, anchor=west},
  wf primary/.style={wf btn, fill=blue!30},
  wf badge/.style={draw=black!35, rounded corners=4pt, fill=white, font=\linespread{1}\fontsize{6}{7}\selectfont\sffamily, inner sep=2pt, anchor=west},
}
% Rozbalovací nabídka od (#1,#2) do (#3,#4) s textem #5.
\newcommand{\wfselect}[5]{%
  \draw[wf field] (#1,#2) rectangle (#3,#4);
  \node[wf t] at ({#1+0.12},{(#2+(#4))/2}) {#5};
  \fill[black!65] ({#3-0.34},{(#2+(#4))/2+0.06}) -- ({#3-0.12},{(#2+(#4))/2+0.06}) -- ({#3-0.23},{(#2+(#4))/2-0.08}) -- cycle;}
% Textové pole od (#1,#2) do (#3,#4) s textem #5.
\newcommand{\wfinput}[5]{%
  \draw[wf field] (#1,#2) rectangle (#3,#4);
  \node[wf t] at ({#1+0.12},{(#2+(#4))/2}) {#5};}
% Zaškrtávací políčko s levým dolním rohem v (#1,#2); #3 = 1 zaškrtnuto; popisek #4.
\newcommand{\wfcheck}[4]{%
  \draw[wf field] (#1,#2) rectangle ({#1+0.26},{#2+0.26});
  \ifnum#3=1 \draw[very thick] ({#1+0.05},{#2+0.13}) -- ({#1+0.11},{#2+0.05}) -- ({#1+0.22},{#2+0.22}); \fi
  \node[wf t] at ({#1+0.36},{#2+0.13}) {#4};}
% Společná hlavička obrazovky široké 16 cm; #1 = 1 zvýrazní Translation, 2 Advisor.
\newcommand{\wfnav}[1]{%
  \fill[wf bar, rounded corners=3pt] (0,0) rectangle (16,-0.7);
  \fill[wf bar] (0,-0.35) rectangle (16,-0.7);
  \node[wf h] at (0.3,-0.35) {ORMConvertor};
  \node[wf t] at (11.0,-0.35) {\ifnum#1=1 \textbf{Translation}\else Translation\fi};
  \node[wf t] at (12.9,-0.35) {\ifnum#1=2 \textbf{Advisor}\else Advisor\fi};
  \node[wf t] at (14.2,-0.35) {Examples};}

\renewcommand{\dateseparator}{.}
\renewcommand{\baselinestretch}{1.5}

\geometry{margin=2cm}

\def\ResearchProject{Research project}
\def\CompanyProject{Company project}
\def\SoftwareProject{Software project}

% METADATA SECTION START
\def\StudentsFullName{Miroslav Štochel}
\def\StudyProgram{Informatika -- Softwarové a datové inženýrství}
\def\ProjectTitle{Rozšíření nástroje ORMorpher}
% Use "\ResearchProject", "\CompanyProject", \SoftwareProject"
\def\ProjectType{\ResearchProject}
\def\SupervisorFullName{Ing. Pavel Koupil, Ph.D.}
\def\ConsultantsFullName{doc. RNDr. Irena Holubová, Ph.D., RNDr. Michal Kopecký, Ph.D.}
% METADATA SECTION END

\title{\ProjectTitle}
\author{\StudentsFullName}

\begin{document}
% Český babel aktivuje spojovník (dělení slov se spojovníkem); TikZ to rozbíjí,
% proto zkratku vypínáme — na sazbu textu to nemá viditelný vliv.
\shorthandoff{-}

\centerline{\Large \textbf{Specification of team software project}}
\centerline{ \textbf{Department of Software Engineering}}
\centerline{ \textbf{Faculty of Mathematics and Physics, Charles University}}
\bigskip
{\noindent\begin{tabular*}{\textwidth}{ >{\raggedleft}m{4cm} l}
 \textbf{Solver:} & \StudentsFullName \\
 \textbf{Study program:} & \StudyProgram \\[1em]
 \textbf{Project title:} & \ProjectTitle \\
 \textbf{Project type:} & \ProjectType \\[1em]
 \textbf{Supervisor:} & \SupervisorFullName \\
 \textbf{Consultants:} & \ConsultantsFullName \\
\end{tabular*}}

\section{Úvod}

V oblasti databázových systémů existuje řada nástrojů a přístupů pro migraci dat, evoluci databázového schématu a propagaci změn do souvisejících datových struktur a dotazů. Výrazně omezenější je však podpora obdobných změn na úrovni aplikačního kódu. Změna databázové technologie, objektově-relačního mapovacího frameworku nebo způsobu perzistence proto obvykle vyžaduje ruční úpravy entit, mapování, dotazů a dalšího kódu datové vrstvy aplikace.

Výzkumná skupina Adaptive Data Management (ADaM) Research Group~\cite{adam-web} se mimo jiné zabývá automatizací migrace a transformace datových a aplikačních artefaktů napříč různými technologiemi. V rámci tohoto výzkumu vzniká nástroj ORMorpher~\cite{ormorpher-web}, jehož původní implementaci vytvořil Milan Abrahám ve své diplomové práci~\cite{abraham-repo} jako proof of concept pro ověření vybraných principů překladu mezi ORM a obdobnými frameworky.

ORMorpher je založen na mezireprezentaci nezávislé na konkrétním frameworku: parser zdrojového frameworku převádí podporované artefakty do mezireprezentace a builder cílového frameworku z ní generuje odpovídající artefakty. Tento přístup omezuje potřebu vytvářet samostatné párové překladače. Současná implementace však podporuje pouze vybrané konstrukce frameworků Dapper, NHibernate a EF~Core a jednotlivé směry překladu nejsou implementovány ve stejném rozsahu. Součástí proof of concept je rovněž Advisor, který experimentálně porovnává vybrané frameworky a umožňuje optimalizační výběr na základě měřených charakteristik.

Cílem projektu je existující proof of concept systematicky rozšířit, nikoli z něj vytvořit produkční nástroj. Projekt doplní složitější mapovací konstrukce, využití metadat z databázového schématu a podporu Java frameworků Hibernate, MyBatis a EclipseLink. Zásadním cílem je tím ověřit použitelnost společné mezireprezentace a překladového přístupu napříč různými programovacími jazyky a technologickými ekosystémy. Součástí projektu je rovněž testovací a běhové prostředí pro automatizované ověřování a experimenty, které umožní posoudit praktickou použitelnost přístupu a připravit srovnání s metodami založenými na velkých jazykových modelech.

Cílem projektu není úplná podpora všech konstrukcí jednotlivých frameworků. Překlad je omezen na explicitně vymezenou společnou podmnožinu konstrukcí, kterou dokáže zachytit mezireprezentace a pro kterou lze ověřit odpovídající chování napříč .NET a Java ekosystémem.

\subsection{Vztah k dalším projektům}

Projekt přímo navazuje na diplomovou práci Milana Abraháma a její prototyp~\cite{abraham-repo}. Dále rozvíjí návrh mezireprezentace, implementaci překladové pipeline a experimentální infrastrukturu původního řešení.

Tematicky souvisí také s nástrojem mm-mapsearch~\cite{mm-mapsearch}, který hledá vhodnou reprezentaci dat s ohledem na zadaný workload, zejména strukturu a charakter dotazů. Zatímco mm-mapsearch řeší volbu reprezentace především na úrovni dat a databázového systému, ORMorpher tuto problematiku doplňuje na úrovni aplikační datové vrstvy transformací entit, mapování a dotazů mezi jednotlivými frameworky.

Další související směr představuje diplomová práce Martina Čorovčáka, která navrhuje univerzální objektový model (UOM) a migraci aplikačního kódu mezi relačními, grafovými a dokumentovými databázovými systémy, včetně mapování ORM, OGM a ODM, s využitím velkých jazykových modelů. Tento přístup je obecnější, ale současně může být nedeterministický a výpočetně i finančně nákladnější. ORMorpher oproti tomu zkoumá deterministický překlad založený na explicitních pravidlech, heuristikách a mezireprezentaci.

S projektem souvisí rovněž diplomová práce Ivony Oboňové zaměřená na migraci dotazů mezi relačními a grafovými databázovými systémy s využitím velkých jazykových modelů. Tyto práce společně pokrývají komplementární části problému: volbu datové reprezentace, transformaci aplikačního mapování a migraci dotazů mezi různými databázovými modely a technologiemi.

Výzkumným přínosem tohoto projektu je především ověření cross-language použitelnosti ORMorpheru rozšířením o Java frameworky, doplnění proof of concept implementace směrem k prakticky použitelnému experimentálnímu nástroji a příprava reprodukovatelného experimentálního prostředí. To poskytne deterministickou baseline pro následné porovnání s LLM-based metodami, včetně variant s RAG nebo agentním přístupem. Výsledky projektu mají sloužit jako podklad pro vědeckou publikaci, přednostně v odborném vědeckém časopise; plnější experimentální porovnání s LLM bude dále rozvíjeno v navazující diplomové práci řešitele (viz kapitolu~\ref{sec:eval}).

\section{Popis projektu}\label{sec:popis}

Nástroj (v~repozitáři ORMConvertor) je backendové řešení v~.NET~10: REST API v~ASP.NET Core s~jednoduchým frontendem. Překlad probíhá ve třech fázích: \emph{parser} zdrojového frameworku převede vstupní artefakty do mezireprezentace, nepovinná fáze \emph{doplnění} dotáhne chybějící mapovací fakta z~katalogu připojené databáze a \emph{builder} cílového frameworku vygeneruje výstupní artefakty. Vedle artefaktů vrací každý převod strukturovanou diagnostiku: záznamy o~ztrátách (fakt, který cíl nevyjádří), konvencích (výstup tvrdí něco, co zdroj neřekl), konfliktech (dva zdroje tvrdí totéž jinak), neúplnostech, selháních a faktech doplněných z~katalogu (výstup tvrdí něco, co dodala databáze, ne zdroj). Dotaz, který by v~cíli vrátil jinou množinu řádků, se zásadně nevydá -- vadný překlad končí záznamem, ne tichou sémantickou změnou. Referenční databází je SQL Server se vzorovou databází WideWorldImporters; celé prostředí včetně testovací sady běží v~Dockeru.

\subsection{Výchozí stav při převzetí}

Snímek stavu k~15.~7.~2026, zdokumentovaný v~repozitáři jako zmražený referenční bod. Prototyp uměl: překlad entit a mapování mezi třemi .NET ORM v~libovolném směru, překlad dotazů jedním směrem (EF~Core LINQ na Dapper SQL), ILP Advisor pro Dapper a EF~Core a REST API s~frontendem. Chybělo: kompozitní primární klíče (model nesl jen jednosloupcový klíč bez pořadí a strategie), vícesloupcové cizí klíče, jakékoli čtení databázového katalogu, priority zdrojů metadat a hlášení konfliktů, validace úplnosti (nepodporované konstrukce se potichu vynechávaly), dotazové parsery pro Dapper a NHibernate, dotazové buildery pro EF~Core a NHibernate, podmínky jiné než plochá konjunkce, záznam běhu a jakákoli podpora javovského ekosystému. Typový model byl vázaný na CLR, takže pro javovské frameworky nepoužitelný.

\subsection{Stav verze 1.2.0}

Práce je rozdělená do tří cílů. \textbf{Cíl~1 -- rozšíření a dotažení tří .NET frameworků -- je hotový} a vydaný jako verze 1.2.0. Nástroj ve verzi 1.2.0 poskytuje:

\begin{itemize}[leftmargin=*]
\item \textbf{Složené identifikátory a vztahy} (F1--F3): mezireprezentace nese klíč jako seřazený seznam částí s~per-part strategií generování a kanonickými parametry generátorů, vztahy s~uspořádanými páry sloupců; 1:1, 1:N i N:M (přes explicitní spojovací entitu) mají uložení a načtení spuštěné proti databázi v~obou cílových frameworcích.
\item \textbf{Doplňování z~databázového katalogu} (F4--F6): jediná komponenta čte metadata z~katalogu SQL Serveru (sloupce, typy, nullabilitu, klíče, unikátní omezení, identity) dávkově; doplní fakta, která zdroj neuvedl, rozpor hlásí konfliktem se zdrojovou hodnotou. Z~entity, která má jen názvy a jazykové typy (typický Dapper vstup), vznikne úplné kompilovatelné mapování včetně vztahů odvozených z~cizích klíčů a detekce spojovacích tabulek.
\item \textbf{Úplná dotazová matice} ve všech devíti směrech mezi třemi frameworky, v~rozsahu projekce, filtrace (podmínkový strom s~AND/OR/NOT a závorkováním), joiny, agregace, řazení, stránkování, poddotazy v~podmínce a množinové operace; NHibernate se čte i z~holého HQL vlastním parserem.
\item \textbf{Validaci a diagnostiku} (F11): brána úplnosti před generováním, šest druhů diagnostických záznamů, záznam běhu s~identifikátorem a verzemi (S6). Vícesouborový vstup a výstup po souborech (část F14).
\item \textbf{Systémové požadavky}: modularita (S1), determinismus (S2), měřený výkon překladu (S3: testovací projekt o~100 entitách a 100 dotazech se přeloží méně než sekundu), kontejnerovou reprodukovatelnost (S5 -- systém i testovací sada jedním příkazem), pozorovatelnost (S6) a základní scénář rozhraní v~pěti krocích (S7).
\end{itemize}

Testovací sada má \textbf{615 testů} (xUnit~v3, 0~přeskočených při připojené databázi) s~měřeným pokrytím 75{,}5~\% řádků; běží v~CI při každém push a pull requestu. Generované artefakty se ověřují na čtyřech stupních: tvar (řetězcové aserce), kompilovatelnost (kompilace Roslynem, validace XML proti XSD), přijetí frameworkem (stavba session factory NHibernate, validace modelu EF~Core, parsování SQL) a běh proti databázi (uložení a načtení entity vygenerovaným kódem). Hranice záruk je vyslovená: Advisor s~benchmarkingem je z~nároku verze vyňatý (nemá testy, měří dva frameworky ze tří), stejně jako dědičnost a komponenty v~NHibernate mapování, cílové dialekty mimo SQL Server a celý javovský ekosystém.

\subsection{Zbývající práce a postup}

\textbf{Cíl~2 -- Java frameworky a cross-language ověření} (F7--F10, F12--F13) je jádrem rozšíření ze zadání: parser a builder pro Hibernate, MyBatis a EclipseLink, překlad mezi .NET a Javou a testovací infrastruktura pro Javu s~diferenčním ověřením výsledků dotazů. Tento cíl má především ověřit, že společná mezireprezentace a překladový přístup nejsou vázané na jediný programovací jazyk nebo technologický ekosystém. Slovníkové předpoklady jsou hotové -- typový model mezireprezentace je jazykově neutrální na jazykové i databázové straně a parametry generátorů klíčů se nesou kanonicky. Prvním krokem cíle je rozhodnutí, kudy javová strana do řešení vstupuje (viz kapitolu~\ref{sec:rizika}). 

\textbf{Cíl~3 -- Advisor, benchmarking a experimenty} (F15, T1--T7) naváže: rozšíření Advisoru o~nové frameworky, izolace spouštění měřeného kódu (S4) a experimentální vyhodnocení podle kapitoly~\ref{sec:eval}. Experimenty mají ověřit praktickou použitelnost rozšířeného proof of conceptu a připravit reprodukovatelnou baseline pro srovnání s LLM-based přístupy.

Postup práce je dokumentovaný přímo v~repozitáři: každá netriviální volba vzniká jako číslované rozhodnutí se zvažovanými variantami a důsledky, aktuální chování popisuje průběžně udržovaný architektonický dokument, zbývající práce jediný seznam otevřených položek a hranici záruk každé vydání vyslovuje explicitně.

\subsection{Funkční požadavky}

\textbf{F1 -- Reprezentace komplexních identifikátorů v~mezireprezentaci.} Entita může mít identifikátor tvořený jednou nebo více vlastnostmi. Mezireprezentace zachová pořadí částí klíče, jejich datové typy, názvy sloupců a strategii generování.

\textbf{F2 -- Překlad komplexních identifikátorů ve všech podporovaných frameworcích.} Složené primární klíče se parsují i generují ve všech podporovaných .NET a Java frameworcích mechanismem daného frameworku (více klíčových sloupců, \texttt{@EmbeddedId}, \texttt{@IdClass}, XML mapování); výstup se zkompiluje a obsahuje všechny části klíče.

\textbf{F3 -- Cizí klíče odkazující na složené identifikátory.} Překlad vztahů s~vícesloupcovým cizím klíčem, včetně 1:1, 1:N a N:M přes spojovací tabulku; každý typ vztahu má zkompilovaný a spuštěný test.

\textbf{F4 -- Načtení metadat z~databázového schématu.} Neobsahuje-li vstupní kód úplná mapovací metadata, systém načte z~připojené databáze alespoň tabulky, sloupce, datové typy, nullabilitu, primární a cizí klíče, unikátní omezení, délku, přesnost a měřítko.

\textbf{F5 -- Sloučení metadat z~více zdrojů.} Informace ze zdrojového kódu, anotací, externích mapovacích souborů a databázového katalogu se slučují do jediné mezireprezentace s~dokumentovanou prioritou zdrojů a hlášením konfliktů (nullabilita, název sloupce, datový typ, primární klíč).

\textbf{F6 -- Generování úplného mapování z~neúplného vstupu.} Z~frameworku s~neúplnou či implicitní mapovací informací (Dapper, MyBatis) vznikne doplněním z~databáze úplné mapování cílového frameworku: z~entity nesoucí jen názvy a jazykové typy vlastností vznikne kompilovatelná cílová entita se správnými názvy tabulek a sloupců, klíči a omezeními.

\textbf{F7 -- Podpora Hibernate.} Parser a generátor pro podporovanou podmnožinu: entity, tabulky a schémata, sloupce, jednoduché a složené klíče, základní vztahy a read-only dotazy alespoň v~JPQL nebo HQL. Testovací sada pokryje všechny implementované kategorie konstrukcí a generované projekty se sestaví.

\textbf{F8 -- Podpora MyBatis.} Parser a generátor pro podporovanou podmnožinu: Java třídy, XML nebo anotované mapování, resultMap, parametry, dynamicky parametrizované read-only dotazy a ručně psané SQL joiny. Testovací sada pokryje implementované kategorie mapování i vícetabulkové dotazy.

\textbf{F9 -- Podpora EclipseLink.} Parser a generátor pro EclipseLink/JPA v~rozsahu společné mezireprezentace: anotované entity, složené klíče, vztahy a JPQL dotazy. Testovací sada pokryje implementované kategorie konstrukcí včetně překladu mezi Hibernate a EclipseLink.

\textbf{F10 -- Cross-language překlad mezi .NET a Javou.} Překlad entity, mapování a podporovaného read-only dotazu mezi implementovanými .NET a Java frameworky v rozsahu společné podporované podmnožiny. Testovací matice bude obsahovat reprezentativní end-to-end scénáře pro směry .NET$\to$Java, Java$\to$.NET a Java$\to$Java.

\textbf{F11 -- Validace cílových artefaktů.} Před předáním výsledku se ověří strukturální úplnost mezireprezentace, omezení cílového frameworku a syntaktická správnost generovaných souborů. Nepodporované konstrukce se nesmí potichu vynechat; každý neúspěšný překlad vrátí strukturovanou diagnostiku (framework, artefakt, chybějící vlastnost, důvod selhání).

\textbf{F12 -- Spustitelná testovací sada pro Java frameworky.} Automatizovaná sada sestaví a spustí vygenerovaný kód Hibernate, MyBatis a EclipseLink nad stejnou databází a daty; je spustitelná jedním příkazem a obsahuje dostatečný počet jednotkových a integračních testů pokrývajících všechny podporované kategorie konstrukcí.

\textbf{F13 -- Diferenční ověření výsledků dotazů.} Zdrojová i přeložená varianta dotazu se spustí a porovnají se normalizované výsledky: pořadí se zohlední jen je-li dané dotazem, práce s~numerickou přesností a hodnotami null je konfigurovatelná. Testovací sada bude obsahovat reprezentativní dvojice dotazů pro podporované kategorie; součástí bude také negativní test se záměrně chybným překladem, který musí být odhalen.

\textbf{F14 -- Překlad celých tříd a dávkových vstupů v~UI.} Vložení či nahrání jedné nebo více celých tříd, případně archivu projektu; volba zdrojového a cílového frameworku; validace a překlad. Rozhraní zobrazí vstup, mezireprezentaci, výstup a diagnostiku po jednotlivých souborech; scénář zpracuje projekt s~alespoň deseti entitami a pěti dotazy bez ručního zadávání každého artefaktu.

\textbf{F15 -- Výběr cílového frameworku a optimalizace (sekundární cíl).} Uživatel může přímo zvolit cílový framework; podle časových možností bude existující Advisor rozšířen také o nové frameworky a umožní doporučení na základě měřeného času, paměti a vah dotazů. Tato funkcionalita je doplňková vůči hlavnímu cíli cross-language překladu a jeho ověření.

\subsection{Nefunkční požadavky}

\textbf{S1 -- Modulární rozšiřitelnost.} Nový framework se přidává novým parserem a builderem přes definovaná rozhraní, bez zásahu do ostatních frameworků a bez párových překladačů.

\textbf{S2 -- Determinismus.} Stejný vstup, konfigurace, schéma a verze nástroje dávají shodné artefakty ve 100~\% běhů.

\textbf{S3 -- Výkon překladu.} Překlad musí být prakticky použitelný pro reprezentativní projektovou testovací sadu; čas samotného překladu a čas načtení databázových metadat se měří a reportují odděleně. Konkrétní výkonnostní výsledky budou součástí vyhodnocení, nikoli pevnou funkční hranicí.

\textbf{S4 -- Izolace a bezpečnost spouštění.} Cizí a generovaný kód v~izolovaném prostředí s~limity; žádné přihlašovací údaje v~artefaktech a lozích. 

\textbf{S5 -- Přenositelné a reprodukovatelné prostředí.} Systém, databáze, testovací projekty a experimentální pipeline v~kontejnerové konfiguraci; reprodukce testů jedním příkazem.

\textbf{S6 -- Pozorovatelnost a auditovatelnost.} Každý překlad má identifikátor běhu a strojově čitelný záznam: zdroje metadat, varování, chyby, verze frameworků i nástroje.

\textbf{S7 -- Uživatelská přívětivost.} Validace před spuštěním, průběžný stav, chyby na úrovni souboru a řádku, stažení výstupu; základní scénář nejvýše v~pěti krocích.

\subsection{Experimentální požadavky}

Experimentální část je rozdělena na požadavky tohoto projektu a navazující experimenty, které budou plně rozvedeny v diplomové práci řešitele. Projekt musí připravit reprodukovatelné testovací prostředí, základní cross-language experimenty a měřicí infrastrukturu; rozsáhlejší LLM experimenty nejsou podmínkou dokončení tohoto projektu.

\textbf{T1 -- Reálná případová studie.} Projekt připraví alespoň jednu reprezentativní případovou studii nebo testovací sadu s netriviálními entitami, vztahy a read-only dotazy. Rozšíření na více reálných open-source aplikací je předmětem navazující diplomové práce.

\textbf{T2 -- Matice překladů.} Projekt ověří směry .NET$\to$Java, Java$\to$.NET a Java$\to$Java na reprezentativních kategoriích dotazů a mapovacích konstrukcí, zejména projekci, filtraci, joinu, agregaci, stránkování, poddotazech a množinových operacích.

\textbf{T3 -- Metriky korektnosti.} Experimentální infrastruktura bude měřit alespoň podíl parsovatelných, kompilovatelných a spustitelných výstupů, funkční ekvivalenci a úplnost přenesených podporovaných mapovacích vlastností.

\textbf{T4 -- Příprava LLM baseline.} Projekt připraví testovací data, rozhraní a měřicí infrastrukturu tak, aby bylo možné stejnou sadu vstupů porovnat s překladem pomocí velkého jazykového modelu. Úplné zero-shot a few-shot experimenty jsou předmětem navazující diplomové práce.

\textbf{T5 -- Příprava RAG / agentní varianty.} Projekt připraví podklady a experimentální rozhraní umožňující v navazující práci doplnit dokumentaci cílového frameworku, informace o schématu a relevantní příklady. Samotné rozsáhlé porovnání RAG a agentních variant není podmínkou dokončení projektu.

\textbf{T6 -- Příprava ablace.} Návrh testovacího prostředí musí umožnit v navazující práci samostatně vyhodnotit vliv databázových metadat, validace, opravného cyklu a případně RAG. Úplná ablační studie není požadavkem tohoto projektu.

\textbf{T7 -- Základní porovnání Advisoru.} Existující Advisor bude podle časových možností rozšířen o nové frameworky a základně porovnán s jednotlivými frameworky a alespoň jednou jednoduchou heuristikou. Jde o sekundární experimentální cíl; nesmí omezit dokončení hlavního cross-language cíle F7--F13.

\subsection{Architektura}

Architekturu popisujeme prvními dvěma úrovněmi modelu C4. Na úrovni kontextu (obr.~\ref{fig:c4l1}) je nástroj jediný systém: uživatel (vývojář migrující aplikaci, nebo výzkumník spouštějící experimenty) s~ním pracuje přes prohlížeč nebo přes REST API, systém čte metadata z~katalogu připojené databáze SQL Server a proti téže databázi měří výkon generovaného kódu. Javovské běhové prostředí je oddělená součást v kontejneru: testovací sada podle F12–F13 si generované artefakty bere z běžící instance přes REST a spouští je proti téže databázi.

\begin{figure}[ht]
\centering
\begin{tikzpicture}[
  node distance=9mm and 26mm,
  every node/.style={font=\small},
  person/.style={rectangle, rounded corners, draw, fill=gray!25, align=center, minimum width=34mm, minimum height=13mm},
  system/.style={rectangle, rounded corners, draw, fill=blue!15, align=center, minimum width=46mm, minimum height=15mm},
  ext/.style={rectangle, rounded corners, draw, fill=gray!10, align=center, minimum width=40mm, minimum height=13mm},
  arr/.style={-latex}
]
\node[person] (user) {Uživatel\\{\scriptsize vývojář / výzkumník}};
\node[system, right=of user] (orm) {\textbf{ORMConvertor}\\{\scriptsize překlad ORM artefaktů,}\\{\scriptsize doporučení frameworku}};
\node[ext, below=of orm] (db) {SQL Server 2022\\{\scriptsize katalog schématu, WideWorldImporters}};
\node[ext, right=of orm, dashed] (java) {Javová testovací sada (JVM)\\{\scriptsize plánováno (F12--F13)}};
\draw[arr] (user) -- node[above]{\scriptsize prohlížeč / REST} (orm);
\draw[arr] (orm) -- node[right]{\scriptsize čtení metadat, měření} (db);
\draw[arr, dashed] (java) -- node[above]{\scriptsize překlady k ověření} (orm);
\draw[arr, dashed] (java.south) |- node[below right]{\scriptsize JDBC} (db.east);
\end{tikzpicture}
\caption{Úroveň kontextu (C4, L1).}
\label{fig:c4l1}
\end{figure}

Na úrovni kontejnerů (obr.~\ref{fig:c4l2}) je systém REST API (ASP.NET Core), překladová pipeline, čtení databázového katalogu a Advisor s~benchmarkingem. Pipeline jsou vzory Adapter (wrappery frameworků), Visitor (generování dotazů z~instrukcí), Builder a Factory; orchestrace pracuje s~wrappery výhradně přes rozhraní. Nosné projekty jsou:

\begin{itemize}[leftmargin=*,itemsep=0pt]
\item \texttt{Model} -- mezireprezentace: entity a mapování (\texttt{EntityMap}, seřazený složený klíč, vztahy s~páry sloupců, unikátní omezení) a dotazové instrukce s~podmínkovým stromem;
\item \texttt{AbstractWrappers} -- rozhraní parserů, abstraktní buildery se šablonovou metodou, deskriptory cílových frameworků (co framework vyžaduje a co nevyjádří) a typy diagnostických záznamů;
\item \texttt{DapperWrappers}, \texttt{EFCoreWrappers}, \texttt{NHibernateWrappers} -- po jednom projektu na framework; sdílené čtení C\# entit a LINQ dotazů nesou projekty \texttt{CSharpEntityParsing} a \texttt{LinqParsing};
\item \texttt{DatabaseCatalog} -- jediné místo čtení databázových metadat a fáze doplnění; wrappery na něm nezávisejí;
\item \texttt{OrmConvertor} -- orchestrace překladu, \texttt{ORMConvertorAPI} -- REST API a frontend, \texttt{Tests} -- testovací sada;
\item \texttt{Advisor} (ILP model v~C nad knihovnou GLPK, volaný přes P/Invoke) a \texttt{AdvisorBenchmarking} (kompilace generovaného kódu Roslynem a měření běhu a paměti).
\end{itemize}

\begin{figure}[ht]
\centering
\begin{tikzpicture}[
  node distance=7mm and 10mm,
  every node/.style={font=\small},
  person/.style={rectangle, rounded corners, draw, fill=gray!25, align=center, minimum width=32mm, minimum height=11mm},
  cont/.style={rectangle, rounded corners, draw, fill=blue!12, align=center, minimum width=40mm, minimum height=12mm},
  ext/.style={rectangle, rounded corners, draw, fill=gray!10, align=center, minimum width=36mm, minimum height=12mm},
  arr/.style={-latex}
]
\node[person] (user) {Uživatel};
\node[cont, below=of user] (api) {REST API + frontend\\{\scriptsize ASP.NET Core, statické stránky}};
\node[cont, below left=9mm and -6mm of api] (pipe) {Překladová pipeline\\{\scriptsize parser $\to$ mezireprezentace}\\{\scriptsize $\to$ doplnění $\to$ builder}};
\node[cont, below right=9mm and -6mm of api] (adv) {Advisor + benchmarking\\{\scriptsize ILP (GLPK), Roslyn,}\\{\scriptsize měření času a paměti}};
\node[cont, below=of pipe] (cat) {DatabaseCatalog\\{\scriptsize dávkové čtení metadat}};
\node[ext, below=of adv] (db) {SQL Server 2022};
\node[draw, rounded corners, dashed, inner sep=3mm, fit=(api)(pipe)(adv)(cat), label={[font=\scriptsize]below left:{ORMConvertor (Docker)}}] {};
\draw[arr] (user) -- (api);
\draw[arr] (api) -- (pipe);
\draw[arr] (api) -- (adv);
\draw[arr] (pipe) -- (cat);
\draw[arr] (adv) -- node[above=1pt, align=center, font=\scriptsize]{překlady\\k~měření} (pipe);
\draw[arr] (cat) -- (db);
\draw[arr] (adv) -- (db);
\end{tikzpicture}
\caption{Úroveň kontejnerů (C4, L2). Javovské kontejnery (F12--F13) přibudou v~cíli~2.}
\label{fig:c4l2}
\end{figure}

Pro každý javovský framework vznikne projekt s novým wrapperem. Provedenou analýzou a popsaným rozhodnutím může implementace rozšířit sdílenou vrstvu.

\subsection{Uživatelské rozhraní}\label{sec:ui}

Rozhraní je záměrně jednoduché: statické stránky se společnou hlavičkou -- rozcestník, překlad, Advisor a výkladová stránka se živými příklady. Základní scénář (S7) je jediná stránka s~pěti očíslovanými kroky viditelnými najednou, takže se uživatel při opravě chyby vrací ke vstupu bez přepínání obrazovek. Makety na obr.~\ref{fig:ui-input}--\ref{fig:ui-advisor} ukazují schematickou podobu obrazovek; obsah polí je ilustrační.

\textbf{Vstup překladu} (obr.~\ref{fig:ui-input}). Kroky 1 a 2 volí zdrojový a cílový framework ze všech podporovaných; nabídka nese verzi z~deskriptoru frameworku, tutéž, kterou pak uvádí záznam běhu (S6). Krok 3 skládá dávkový vstup podle F14 jako seznam pojmenovaných jednotek: jednotka má jméno, typ obsahu a obsah a vzniká nahráním jednoho či více celých souborů (typ se předvyplní podle přípony), přidáním prázdné jednotky a vepsáním, nebo naplněním vzorky. Nabídka typů je omezená na to, co zdrojový framework čte. Krok 4 nejprve validuje na klientovi -- prázdný seznam, prázdná jednotka, typ mimo nabídku, nesprávně utvořené XML s~číslem řádku -- a chybu ukáže u~dotčené jednotky; syntaktické chyby zdrojového kódu hlásí server a rozhraní zobrazí jeho hlášku doslova, u~SQL s~řádkem a sloupcem. Po dobu volání je tlačítko zablokované a průběžný stav se hlásí pod ním (S7).

\begin{figure}[htbp]
\centering
\begin{tikzpicture}[x=1cm,y=1cm]
\draw[wf frame] (0,0) rectangle (16,-11.55);
\wfnav{1}
\node[wf h, font=\linespread{1}\fontsize{11}{12}\selectfont\sffamily\bfseries] at (0.4,-1.1) {Translation};
% krok 1
\node[wf h] at (0.4,-1.7) {1. Source framework};
\wfselect{0.4}{-1.95}{6.4}{-2.45}{Entity Framework Core 10.0.10}
% krok 2
\node[wf h] at (0.4,-2.9) {2. Target framework};
\wfselect{0.4}{-3.15}{6.4}{-3.65}{NHibernate 5.7.0}
% krok 3
\node[wf h] at (0.4,-4.55) {3. Input files};
\node[wf primary] at (0.4,-5.0) {Add files\dots};
\node[wf btn] at (1.95,-5.0) {Add empty unit};
\node[wf btn] at (3.95,-5.0) {Load sample set};
% jednotka 1
\draw[wf card] (0.4,-5.4) rectangle (15.6,-7.45);
\wfinput{0.55}{-5.53}{4.6}{-5.98}{Customer.cs}
\wfselect{4.8}{-5.53}{8.0}{-5.98}{C\# entity}
\draw[wf pre] (0.55,-6.12) rectangle (15.45,-7.33);
\node[wf code] at (0.65,-6.2) {[Table("Customers", Schema = "Sales")]\\public class Customer\\\{\\~~~~[Key] public int CustomerID \{ get; set; \}\\~~~~[MaxLength(100)] public string CustomerName \{ get; set; \} = null!;};
% jednotka 2
\draw[wf card] (0.4,-7.6) rectangle (15.6,-9.1);
\wfinput{0.55}{-7.73}{4.6}{-8.18}{CustomerByCity.cs}
\wfselect{4.8}{-7.73}{8.0}{-8.18}{C\# LINQ query}
\draw[wf pre] (0.55,-8.32) rectangle (15.45,-8.98);
\node[wf code] at (0.65,-8.4) {var customers = context.Customers.Where(c => c.DeliveryCityID == cityId)\\~~~~.OrderBy(c => c.CustomerName).Select(c => new \{ c.CustomerID, c.CustomerName \});};
% jednotka 3 (chyba validace)
\draw[wf card] (0.4,-9.25) rectangle (15.6,-10.55);
\wfinput{0.55}{-9.38}{4.6}{-9.83}{Order.cs}
\wfselect{4.8}{-9.38}{8.0}{-9.83}{C\# entity}
\draw[wf pre, draw=red!65!black] (0.55,-9.97) rectangle (15.45,-10.27);
\node[wf err] at (0.55,-10.42) {The unit has no content.};
% krok 4
\node[wf h] at (0.4,-10.85) {4. Convert};
\node[wf primary] at (0.4,-11.2) {Convert};
% otevřená nabídka cílů (kreslí se naposled, leží přes obsah)
\draw[wf popup] (7.3,-3.15) rectangle (11.6,-5.37);
\fill[blue!14] (7.32,-3.53) rectangle (11.58,-3.88);
\node[wf t] at (7.45,-3.35) {Dapper 2.1.79};
\node[wf t] at (7.45,-3.705) {NHibernate 5.7.0};
\node[wf t] at (7.45,-4.06) {Entity Framework Core 10.0.10};
\node[wf t] at (7.45,-4.415) {Hibernate ORM 7.4.5};
\node[wf t] at (7.45,-4.77) {MyBatis 3.5.19};
\node[wf t] at (7.45,-5.125) {EclipseLink 5.0.0};
\end{tikzpicture}
\caption{Překladová obrazovka, kroky 1--4: volba frameworků, dávkový vstup po pojmenovaných jednotkách a validace před odesláním.}
\label{fig:ui-input}
\end{figure}

\textbf{Výsledek překladu} (obr.~\ref{fig:ui-result}). Hlavička nese záznam běhu (S6): zdroj $\to$ cíl s~verzemi, stav připojení ke katalogu (nenakonfigurováno, nepoužito, přečteno s~časem čtení, nedosažitelné), verzi nástroje a identifikátor běhu. Artefakty stojí po souborech v~panelech se zvýrazněnou syntaxí; kompletní výstup balí server do ZIP (S7). Diagnostické záznamy (F11) stojí pod výstupem: souhrn počtů podle druhu slouží zároveň jako filtr a tabulka nese druh, entitu, vlastnost, vstupní jednotku, cílový artefakt, předmět a důvod; záznamy druhu \emph{Failure} stojí první, protože znamenají, že artefakty entity chybí. Jména souborů výstupu odvozuje klient z~obsahu artefaktu.

\begin{figure}[htbp]
\centering
\begin{tikzpicture}[x=1cm,y=1cm]
\draw[wf frame] (0,0) rectangle (16,-8.3);
\node[wf h] at (0.4,-0.35) {5. Result};
\draw[wf card] (0.4,-0.6) rectangle (15.6,-1.15);
\node[wf t] at (0.55,-0.875) {Entity Framework Core 10.0.10 $\to$ NHibernate 5.7.0};
\node[wf badge, fill=green!12] at (6.7,-0.875) {catalog read (38 ms)};
\node[wf t] at (9.6,-0.875) {ORMConvertor \texttt{3.0.0}};
\node[wf t] at (12.2,-0.875) {run \texttt{2f8c1e7a\dots}};
\node[wf btn] at (0.4,-1.5) {Download all as ZIP};
% artefakt A
\draw[wf card] (0.4,-1.9) rectangle (7.9,-4.3);
\node[wf h] at (0.55,-2.15) {Customer.cs};
\node[wf tt] at (2.4,-2.15) {C\# entity};
\draw[wf pre] (0.55,-2.4) rectangle (7.75,-4.17);
\node[wf code] at (0.65,-2.47) {public class Customer\\\{\\~~~~public virtual int CustomerID \{ get; set; \}\\~~~~public virtual string CustomerName \{ get; set; \} = null!;\\~~~~public virtual City DeliveryCity \{ get; set; \} = null!;\\\}};
% artefakt B
\draw[wf card] (8.1,-1.9) rectangle (15.6,-4.3);
\node[wf h] at (8.25,-2.15) {Customer.hbm.xml};
\node[wf tt] at (11.0,-2.15) {XML mapping};
\draw[wf pre] (8.25,-2.4) rectangle (15.45,-4.17);
\node[wf code] at (8.35,-2.47) {<class name="Customer" table="Customers" schema="Sales">\\~~<id name="CustomerID" column="CustomerID">\\~~~~<generator class="identity" />\\~~</id>\\~~<property name="CustomerName" length="100" not-null="true" />\\~~<many-to-one name="DeliveryCity" column="DeliveryCityID" />\\</class>};
% artefakt C
\draw[wf card] (0.4,-4.5) rectangle (7.9,-6.3);
\node[wf h] at (0.55,-4.75) {CustomerByCity.hql};
\node[wf tt] at (3.5,-4.75) {HQL query};
\draw[wf pre] (0.55,-5.0) rectangle (7.75,-6.17);
\node[wf code] at (0.65,-5.07) {select c.CustomerID, c.CustomerName\\from Customer c\\where c.DeliveryCity.CityID = :cityId\\order by c.CustomerName};
% mezireprezentace
\draw[wf card] (8.1,-4.5) rectangle (15.6,-6.3);
\node[wf h] at (8.25,-4.75) {Intermediate representation};
\node[wf tt] at (12.6,-4.75) {Customer};
\draw[wf pre] (8.25,-5.0) rectangle (15.45,-6.17);
\node[wf code] at (8.35,-5.07) {EntityMap Customer -> Sales.Customers\\~~key       CustomerID : int, generated (identity)\\~~property  CustomerName : string(100), not null\\~~relation  DeliveryCity -> City, N:1, owning\\~~~~~~~~~~~~DeliveryCityID -> CityID};
% záznamy
\node[wf h] at (0.4,-6.6) {Records};
\node[wf badge, fill=yellow!20] at (0.4,-6.95) {Convention $\times$ 1};
\node[wf badge, fill=green!12] at (2.1,-6.95) {Supplied $\times$ 1};
\node[wf tab] at (0.4,-7.25) {\setlength{\tabcolsep}{3pt}\begin{tabular}{@{}lllllll@{}}
\textbf{Kind} & \textbf{Entity} & \textbf{Property} & \textbf{Unit} & \textbf{Artifact} & \textbf{Subject} & \textbf{Reason}\\ \hline
Convention & Customer & CustomerID & Customer.cs & \underline{Customer.hbm.xml} & key generation & no strategy stated; identity assumed for an integer key\\
Supplied & Customer & CustomerName & Customer.cs & \underline{Customer.hbm.xml} & nullability & not-null read from the catalog (Sales.Customers)\\
\end{tabular}};
\end{tikzpicture}
\caption{Překladová obrazovka, krok 5: záznam běhu se stavem katalogu, artefakty po souborech, mezireprezentace a diagnostické záznamy po jednotkách.}
\label{fig:ui-result}
\end{figure}

\textbf{Advisor} (obr.~\ref{fig:ui-advisor}) sdílí tvar s~překladem: zdrojový framework, entitní jednotky, dotazy s~vahami (váha je relativní četnost dotazu), omezení -- limit paměti, strop počtu vybíraných frameworků a podmnožina kandidátů -- a spuštění, které pro každý dotaz a každého kandidáta překládá, kompiluje a měří, a proto trvá minuty. Výsledek ukazuje srovnání: pro každý dotaz všechny změřené kandidáty seřazené podle času, s~vyznačeným přiřazením, poměrem k~nejrychlejšímu a součty za přiřazené řádky. Přiřazený kandidát nemusí být nejrychlejší, protože limit paměti a strop počtu frameworků jsou součástí modelu -- v~ukázce je strop jeden framework, takže druhý dotaz dostal Dapper, ač je v~EF~Core rychlejší. Vedle srovnání stojí tabulka variant podle F15: pro každého kandidáta varianta \uv{vše v~jednom frameworku}, heuristika volící pro každý dotaz nejrychlejšího kandidáta a ILP nad shodným workloadem.

\begin{figure}[htbp]
\centering
\begin{tikzpicture}[x=1cm,y=1cm]
\draw[wf frame] (0,0) rectangle (16,-13.7);
\wfnav{2}
\node[wf h, font=\linespread{1}\fontsize{11}{12}\selectfont\sffamily\bfseries] at (0.4,-1.1) {Advisor};
% krok 1
\node[wf h] at (0.4,-1.7) {1. Source framework};
\wfselect{0.4}{-1.95}{6.4}{-2.45}{Entity Framework Core 10.0.10}
% krok 2
\node[wf h] at (0.4,-2.9) {2. Entities};
\node[wf btn] at (0.4,-3.25) {Load samples};
\draw[wf card] (0.4,-3.6) rectangle (7.9,-4.9);
\wfinput{0.55}{-3.73}{3.6}{-4.18}{Customer.cs}
\wfselect{3.75}{-3.73}{5.9}{-4.18}{C\# entity}
\draw[wf pre] (0.55,-4.32) rectangle (7.75,-4.78);
\node[wf code] at (0.65,-4.39) {[Table("Customers", Schema = "Sales")]\\public class Customer \{ \dots \}};
\draw[wf card] (8.1,-3.6) rectangle (15.6,-4.9);
\wfinput{8.25}{-3.73}{11.3}{-4.18}{Order.cs}
\wfselect{11.45}{-3.73}{13.6}{-4.18}{C\# entity}
\draw[wf pre] (8.25,-4.32) rectangle (15.45,-4.78);
\node[wf code] at (8.35,-4.39) {[Table("Orders", Schema = "Sales")]\\public class Order \{ \dots \}};
% krok 3
\node[wf h] at (0.4,-5.25) {3. Queries and weights};
\node[wf btn] at (0.4,-5.6) {Add query};
\draw[wf card] (0.4,-5.95) rectangle (15.6,-6.95);
\wfinput{0.55}{-6.08}{4.0}{-6.53}{CustomerByCity.cs}
\node[wf t] at (4.2,-6.305) {weight};
\wfinput{5.0}{-6.08}{6.0}{-6.53}{10}
\draw[wf pre] (0.55,-6.63) rectangle (15.45,-6.88);
\node[wf code] at (0.65,-6.68) {context.Customers.Where(c => c.DeliveryCityID == cityId).Select(c => c.CustomerName)};
\draw[wf card] (0.4,-7.1) rectangle (15.6,-8.1);
\wfinput{0.55}{-7.23}{4.0}{-7.68}{OrdersOfCustomer.cs}
\node[wf t] at (4.2,-7.455) {weight};
\wfinput{5.0}{-7.23}{6.0}{-7.68}{3}
\draw[wf pre] (0.55,-7.78) rectangle (15.45,-8.03);
\node[wf code] at (0.65,-7.83) {context.Orders.Where(o => o.CustomerID == id).OrderByDescending(o => o.OrderDate).Take(20)};
% krok 4
\node[wf h] at (0.4,-8.45) {4. Constraints};
\node[wf tt] at (0.4,-8.72) {Memory limit (KB)};
\wfinput{0.4}{-8.87}{3.2}{-9.3}{0}
\node[wf tt] at (4.6,-8.72) {Maximum frameworks to select};
\wfinput{4.6}{-8.87}{7.4}{-9.3}{1}
\node[wf tt] at (8.2,-8.72) {Candidate frameworks};
\wfcheck{8.2}{-9.2}{1}{Dapper}
\wfcheck{10.4}{-9.2}{0}{NHibernate}
\wfcheck{12.3}{-9.2}{1}{Entity Framework Core}
\wfcheck{8.2}{-9.65}{1}{Hibernate}
\wfcheck{10.4}{-9.65}{0}{MyBatis}
\wfcheck{12.3}{-9.65}{0}{EclipseLink}
% krok 5
\node[wf h] at (0.4,-10.05) {5. Run};
\node[wf primary] at (0.4,-10.4) {Run the advisor};
% výsledek
\node[wf h] at (0.4,-10.85) {Result};
\node[wf t] at (0.4,-11.15) {Selected: \textbf{Dapper} (1 framework allowed for the whole workload)};
\node[wf tab] at (0.4,-11.45) {\setlength{\tabcolsep}{3pt}\begin{tabular}{@{}llrrrl@{}}
\textbf{Query} & \textbf{Framework} & \textbf{Mean time (ms)} & \textbf{Allocated (KB)} & \textbf{Relative} & \\ \hline
CustomerByCity & Dapper & 0.41 & 12 & 1.00 & assigned\\
 & EF Core & 0.63 & 41 & 1.54 & \\
 & Hibernate & 0.71 & 64 & 1.73 & \\
OrdersOfCustomer & EF Core & 1.20 & 96 & 1.00 & \\
 & Dapper & 1.25 & 30 & 1.04 & assigned\\
 & Hibernate & 1.42 & 118 & 1.18 & \\ \hline
\textbf{Total (assigned)} & & \textbf{1.66} & \textbf{42} & & \\
\end{tabular}};
\node[wf h] at (9.7,-10.85) {Variants of the same workload};
\node[wf tab] at (9.7,-11.45) {\setlength{\tabcolsep}{3pt}\begin{tabular}{@{}lrr@{}}
\textbf{Variant} & \textbf{Time (ms)} & \textbf{Memory (KB)}\\ \hline
all Dapper & 1.66 & 42\\
all EF Core & 1.83 & 137\\
all Hibernate & 2.13 & 182\\
heuristic (fastest per query) & 1.61 & 108\\
ILP (selected) & 1.66 & 42\\
\end{tabular}};
\end{tikzpicture}
\caption{Advisor: vstup s~vahami dotazů a omezeními, výsledek jako srovnání všech změřených kandidátů s~vyznačeným přiřazením a tabulka variant podle F15.}
\label{fig:ui-advisor}
\end{figure}

\section{Platformy a technologie}

Nástroj je .NET~10 solution; všechny verze závislostí jsou zafixované a spravované centrálně na jediném místě repozitáře.

\begin{itemize}[leftmargin=*,itemsep=0pt]
\item \textbf{Backend:} C\#, ASP.NET Core (REST API, OpenAPI dokument), Roslyn (čtení C\# entit a LINQ, kompilační stupeň ověření), \texttt{Microsoft.SqlServer.TransactSql.ScriptDom} (parsování T-SQL). Platforma, REST API i Roslyn jsou zděděné z~prototypu, který běžel na .NET~8. Verzi jsme změnili na aktuální .NET~10. ScriptDom čte holé SQL Dapperu a mapperů MyBatis.
\item \textbf{Překládané .NET frameworky:} Dapper 2.1.79, NHibernate 5.7.0, Entity Framework Core 10.0.10. Frameworky jsou zděděné z~prototypu, zafixovali jsme jejich verze odpovídající .NET~10.
\item \textbf{Java ekosystém (cíl~2, verze zafixované pro projekt):} JDK~25 (LTS), Jakarta Persistence 3.2, Hibernate ORM 7.4.5, EclipseLink 5.0.0, MyBatis 3.5.19, ovladač \texttt{mssql-jdbc}; sestavení Maven; JVM žije jen v kontejneru testovací sady a na CI runneru. Frameworky předpokládá zadání, verze JDK~25 je jako LTS protějšek .NET~10, u~Hibernate a EclipseLink zvolené implementace stejné verze Jakarta Persistence a MyBatis 3.5.19 jako poslední vydání řady 3.5.
\item \textbf{Databáze:} SQL Server 2022, vzorová databáze WideWorldImporters; metadata z~pohledů \texttt{sys.*}. Databáze i vzorová data jsou zděděné z~prototypu, kde nad nimi měří Advisor.
\item \textbf{Optimalizace:} GLPK (GNU Linear Programming Kit) volaný z~C~wrapperu přes P/Invoke. Zděděné z~prototypu.
\item \textbf{Frontend:} ručně psané statické HTML, nativní ES moduly a CSS, vendorované Pico~CSS a highlight.js. Prototyp měl frontend v~Angularu s~commitnutým zastaralým buildem. Statické stránky nepotřebují build a nasazení nezávisí na npm ekosystému. Pokud se při používání ukáže nedostatečným, přebuduje se frontend do některého JavaScriptového frameworku.
\item \textbf{Testy a infrastruktura:} xUnit~v3, GitHub Actions (CI), Docker a docker~compose (aplikace, databáze i testovací sada jedním příkazem). Testovací sada, CI i kontejnerová konfigurace aplikace s~databází jsou zděděné z~prototypu. xUnit ve verzi~3 umožňuje dynamické přeskakování testů závislých na databázi s~uvedeným důvodem.
\end{itemize}

Seznam není definitivní -- konkrétní volby javovské strany (způsob parsování Java kódu, sestavovací nástroj, tvar testovacího harnessu) určí analýza a samostatná rozhodnutí.

\section{Vyhodnocení}\label{sec:eval}

Vyhodnocení má dvě roviny: průběžnou kontrolu kvality nástroje a experimenty, které ověřují hlavní cíle projektu.

\textbf{Kvalita nástroje.} Generované artefakty se ověřují ve více stupních: (1)~tvar výstupu a strukturální úplnost, (2)~kompilovatelnost nebo syntaktická validita, (3)~přijetí cílovým frameworkem a (4)~u vybraných scénářů běh proti databázi. Pro Java frameworky vznikne analogická testovací infrastruktura v rozsahu podporované podmnožiny. Požadavky F12--F13 k tomu přidávají spustitelnou sadu a diferenční ověření výsledků dotazů nad stejnými daty.

\textbf{Experimenty projektu.} Jádrem experimentálního vyhodnocení je cross-language ověření: překlady .NET$\to$Java, Java$\to$.NET a Java$\to$Java, jejich kompilovatelnost, spustitelnost, funkční ekvivalence a úplnost přenesených podporovaných mapovacích vlastností. Tyto experimenty mají ověřit, že mezireprezentace a překladový přístup nejsou vázány na jediný programovací jazyk a že rozšířený proof of concept je prakticky použitelný pro reprezentativní scénáře.

\textbf{Navazující LLM experimenty.} Projekt připraví testovací data a měřicí infrastrukturu tak, aby bylo možné deterministický pravidlový překlad porovnat se stejnými vstupy zpracovanými pomocí velkých jazykových modelů. Plné zero-shot, few-shot, RAG, agentní a ablační experimenty budou rozvedeny v navazující diplomové práci řešitele; nejsou podmínkou dokončení tohoto projektu.

\textbf{Advisor a optimalizace.} Existující Advisor a ILP optimalizace představují sekundární experimentální větev. Pokud to dovolí časový rozsah projektu, budou základně rozšířeny o nové frameworky a porovnány s jednoduchými baseline strategiemi. Dokončení hlavního cross-language cíle má před touto částí prioritu.

Výsledky experimentů provedených v~rámci projektu budou předloženy u~obhajoby spolu s~naměřenými ukazateli kvality nástroje. Zvláštní důraz bude kladen na ověření cross-language použitelnosti přístupu rozšířením o Java frameworky, praktickou použitelnost rozšířené proof of concept implementace a reprodukovatelné experimentální ověření. Tyto výsledky mají sloužit jako podklad pro vědeckou publikaci, přednostně v odborném vědeckém časopise, a jako deterministická baseline pro následné plnější porovnání s LLM-based metodami.

\section{Rizika}\label{sec:rizika}

\textbf{Vstup javovské strany do řešení.} Zásadní nejistota cíle~2: buď javovské wrappery čtou a píší javovský kód vlastními parsery v~C\# (jednodušší nasazení, ale žádné ověření kompilací a přijetím frameworku), nebo k~řešení přibude javovská komponenta s~JVM (plné ověřovací stupně a diferenční běhy, ale druhý runtime v~kontejnerové konfiguraci). Riziko se řeší jako první rozhodnutí cíle~2, s~vědomím, že diferenční ověření (F13) JVM s~JDBC připojením vyžaduje v~každém případě -- varianty se liší tím, kde JVM stojí, ne tím, jestli existuje.

\textbf{Parsování javovských jazyků bez referenčních parserů.} Jediný úplný parser HQL/JPQL žije uvnitř příslušného frameworku a dynamické SQL MyBatisu je vlastní XML dialekt. Zmírnění: stejný postup, jaký se osvědčil u~HQL v~cíli~1 -- vlastní parser podmnožiny, kterou mezireprezentace unese, držený round-tripovými testy proti vlastnímu builderu; konstrukce mimo podmnožinu končí definovaným záznamem, ne tichým vynecháním.

\textbf{SQL dialekty.} Ručně psané SQL v~MyBatis mapperech bývá psané pro jiné databáze než SQL Server a stávající T-SQL gramatika ho nepřečte. Zmírnění: rozsah se drží na referenčním SQL Serveru (jediný nárokovaný dialekt); co gramatika nepřečte, se hlásí záznamem o~neúplnosti.

\textbf{Sémantické rozdíly frameworků.} Konstrukce, kterou cíl nevyjádří, nesmí zmizet potichu ani změnit význam. Mechanismus je hotový a ověřený z~cíle~1: deskriptor cílového frameworku deklaruje, co framework nevyjádří, ztráty se hlásí mechanicky a dotaz s~jinou množinou řádků se nevydá; javovské frameworky dostanou vlastní deskriptory. Zbytkové riziko je pracnost analýzy každého frameworku proti zafixované verzi.

\textbf{Metodologie měření pro Javu.} Advisor dnes měří .NET kód kompilovaný Roslynem v~procesu; pro Javu je třeba najít srovnatelné běhové prostředí a vyrovnat se s~rozdíly JIT, startovní režie a správy paměti. Zmírnění: analýza prostředí je samostatný krok harmonogramu; iterační politika měření bude vyslovena a odůvodněna jako součást metodologie, aby čísla byla obhajitelná (T7). S~tím souvisí i dosud nevyřešená izolace spouštění měřeného kódu (S4) -- do doby rozhodnutí platí předpoklad nasazení v~důvěryhodné síti a oblast je výslovně vyňata ze záruk.

\textbf{Pracnost a čas.} Tři nové frameworky v reprezentativním rozsahu společné mezireprezentace představují objemné zadání. Zmírnění: cíle jsou seřazené tak, aby jádro zadání (F7--F10) šlo před experimenty; mezireprezentace i slovníky jsou na Javu připravené z~cíle~1; rozsah experimentů je odstupňovaný (viz kapitolu~\ref{sec:eval}) a plné LLM experimenty nese až diplomová práce. Průběh je průběžně viditelný v~repozitáři, takže případný skluz je vidět včas a lze ho řešit se supervizorem.

\section{Milníky / Výstupy}

Projekt se odevzdává jako veřejný repozitář na GitHubu~\cite{fork-repo} obsahující zdrojové kódy, kontejnerovou konfiguraci (aplikace, databáze i testovací sada jedním příkazem), testovací sadu a dokumentaci: průběžně udržovaný popis architektury, číslovaná návrhová rozhodnutí, seznam otevřených položek a explicitní hranici záruk každého vydání. Průběh prací a podíl řešitele je zřejmý z~commit historie a tagů vydání.

\begin{itemize}[leftmargin=*,itemsep=0pt]
\item \textbf{Milník 1 -- dokončená .NET část:} složené klíče, vztahy, katalogové doplnění, diagnostika, úplná dotazová matice tří .NET frameworků, čtyři stupně ověření.
\item \textbf{Milník 2 -- specifikace a analýza:} tento dokument; rozhodnutí o~vstupu javovské strany a analýza běhového prostředí pro Javu.
\item \textbf{Milník 3 -- Java frameworky a cross-language ověření:} parsery a buildery Hibernate, MyBatis a EclipseLink, překlad mezi .NET a Javou a ověření použitelnosti společné mezireprezentace napříč oběma ekosystémy.
\item \textbf{Milník 4 -- testovací sady a experimenty:} Java testovací infrastruktura s~diferenčním ověřením (F12--F13), zobrazení mezireprezentace, cross-language matice T2 s~metrikami T3 a příprava infrastruktury pro navazující LLM experimenty; základní srovnání Advisoru (T7) podle časových možností.
\item \textbf{Milník 5 -- odevzdání:} cross-language experimentální výsledky, finální dokumentace, odevzdání a příprava obhajoby s~živou ukázkou nástroje.
\end{itemize}

Komisi bude u~obhajoby předvedena běžící instance (docker~compose) na překladu reálných artefaktů mezi .NET a Javou, reprodukce testovací sady jedním příkazem a výsledky experimentů. Dílčí demonstrace supervizorovi a konzultantovi probíhají při uzavírání milníků.

\section{Harmonogram}

\begin{tabular}{ | m{11em} | m{23em} | m{3cm} | m{1cm} | }
   \hline
 \textbf{Milník / výstup} & \textbf{Aktivita} & \textbf{Časový plán} & \textbf{MD} \\
 \hline
 1. Dokončení .NET části & seznámení, audit prototypu, dokumentační systém & červenec & 4 \\
                        & klíče, vztahy, katalog, diagnostika, dotazová matice, ověření; vydání 1.2.0 & červenec--srpen & 9 \\
 \hline
 2. Specifikace a analýza & specifikace; rozhodnutí o~javovské straně, analýza běhového prostředí & září & 5 \\
 \hline
 3. Javovské frameworky & Hibernate (parser + builder) & září & 4 \\
                        & MyBatis & září--říjen & 3 \\
                        & EclipseLink, cross-ecosystem testy & září--říjen & 3 \\
 \hline
 4. Testy a experimenty & javovská testovací sada, diferenční ověření (F12--F13), zobrazení mezireprezentace & říjen--listopad & 5 \\
 \hline
 5. Odevzdání & experimenty (T2, T3, T7), dokumentace, odevzdání & listopad & 5 \\
 \hline
\end{tabular}

\section{Forma spolupráce}

Projekt řeší jediný řešitel v~rámci výzkumné skupiny Adaptive Data Management (ADaM) Research Group. Supervizorem projektu je Pavel Koupil, který je autorem zadání a odborným garantem doménové části, zejména návrhu mezireprezentace, překladových pravidel, návaznosti na související projekty a metodologie experimentů. Současně dohlíží na průběh projektu, softwarově-inženýrskou kvalitu a plnění harmonogramu.

Řešitel se pravidelně (typicky jednou za týden) setkává se supervizorem k diskusi nad provedenými změnami, návrhovými rozhodnutími a dalším plánem vývoje. Návrhová rozhodnutí a průběh implementace jsou průběžně dokumentovány v repozitáři.

Vzhledem k návaznosti na další projekty výzkumné skupiny se počítá také s příležitostnou komunikací s jejich autory. V oblasti univerzálního objektového mapování může řešitel konzultovat zejména s Martinem Čorovčákem otázky související s reprezentací a migrací mapování mezi relačními, grafovými a dokumentovými systémy. Pro synchronizaci s projektem mm-mapsearch a souvisejícími pracemi může podle potřeby komunikovat s Bedřichem Mazourkem nebo Jáchymem Bártíkem. V době této spolupráce již mohou mít tito studenti své projekty dokončené nebo obhájené; jejich role je proto chápána především jako konzultace s autory předchozích řešení a předávání zkušeností.

Projekt zároveň předpokládá spolupráci s novými členy týmu, kteří se budou zapojovat do navazujících témat datové migrace, mapování a experimentálního vyhodnocení. Podle potřeby budou jednotlivé otázky konzultovány také s Irenou Holubovou nebo Michalem Kopeckým, zejména v souvislosti s vlastnostmi relačního datového modelu, databázových schémat a sémantikou relačních dat.

\begin{thebibliography}{9}


\bibitem{adam-web}
ADAPTIVE DATA MANAGEMENT RESEARCH GROUP. \emph{Adaptive Data Management Research Group}. Dostupné z: \url{https://www.cnaip.cz/en/entity/adaptive-data-management-research-group}

\bibitem{mm-mapsearch}
KOUPIL, Pavel et al. \emph{mm-mapsearch}. Stránka nástroje. Dostupné z: \url{https://www.ksi.mff.cuni.cz/~koupil/software/mm-mapsearch/}

\bibitem{ormorpher-web}
KOUPIL, Pavel. \emph{ORMorpher}. Stránka nástroje. Dostupné z: \url{https://www.ksi.mff.cuni.cz/~koupil/ormorpher/index.html}

\bibitem{abraham-repo}
ABRAHÁM, Milan. \emph{orm-convertor} (ORMorpher). Repozitář prototypu a text diplomové práce. Dostupné z: \url{https://github.com/milan252525/orm-convertor}

\bibitem{fork-repo}
ŠTOCHEL, Miroslav. \emph{orm-convertor} (ORMConvertor). Repozitář projektu. Dostupné z: \url{https://github.com/MiraStochel/orm-convertor}

\end{thebibliography}

\end{document}